\subsubsection{Additional statistical, recurrent, and attention baselines}
\label{sec:additional_baselines}

The additional benchmark retains persistence, annual seasonal na\"ive,
climatology where previously evaluated, XGBoost, and LSTM. It adds ARIMA,
GRU, a full-attention Transformer encoder--decoder, and a compact Informer.
All additions use the same ten sites, partitions, forecast-origin keys,
targets, normalization, horizons, and physical post-processing. Original
forecasts are read without alteration; equality of origin--lead keys and
targets is verified before training and again during consolidation.

The GRU \cite{cho2014gru} replaces the canonical LSTM cell while retaining
a one-layer, 32-unit encoder, an eight-dimensional site embedding, and a
linear direct-output head. The Transformer \cite{vaswani2017attention}
uses complete multihead attention, two encoder layers, one decoder layer,
model width 16, feed-forward width 32, two heads, ReLU, and dropout 0.1.
Linear GHI projections and sinusoidal relative positions replace the
original language-model input representation. A learned eight-dimensional
site embedding is projected to the model width. The decoder receives the
last half of the historical context followed by zeros for the entire
forecast horizon; neither training nor inference supplies future targets
or meteorological covariates.

Informer \cite{zhou2021informer} uses the same compact widths and depth,
but replaces encoder and causal decoder self-attention with ProbSparse
attention (factor 5), inserts convolution--batch-normalization--ELU--pooling
distillation between encoder layers, and produces the complete forecast
in one decoder pass. Cross-attention remains complete. Evaluation uses
fixed sampled-key indices independent of batch membership; training uses
the registered random seed. At short sequence lengths, all queries can be
selected, as implied by the capped ProbSparse rule. These are explicit
GHI-only adaptations, not reproductions of the original papers' calendar
embeddings, parameter budgets, or reported benchmark scores.

All three added networks use Adam, learning rate $10^{-3}$, weight decay
$10^{-5}$, and one CPU thread, matching the canonical optimization policy.
Hourly and daily tasks use batches of 128, at most 30 selection epochs,
and patience 5. Monthly tasks use batches of 32, at most 300 epochs, and
patience 30. Early stopping minimizes normalized validation MSE in 2023.
A freshly initialized model with the same seed is then refitted for the
selected number of epochs on the pre-test training and validation data.
Hourly runs use seed 42; daily and monthly runs use seeds 11, 23, 42, 67,
and 89, with ensemble formation matching Eq.~\eqref{eq:ensemble_forecast}.

ARIMA is fitted separately by site using statsmodels 0.14.6
\cite{statsmodels2026arima}. The candidate $(p,d,q)$ orders are
$(1,0,0)$, $(2,0,0)$, $(2,0,2)$, $(1,1,0)$, $(0,1,1)$, $(1,1,1)$, and
$(2,1,2)$. Parameters are estimated from unique observations in the
admissible training-target date span. The order minimizes the site's
normalized validation MSE before physical post-processing. The selected
order is refitted on the corresponding pre-2024 target span. At each
origin, fixed parameters are applied to the same bounded historical window
used by the neural predictors; the future vector is generated recursively
without target feedback. Thus, the forecast targets are matched even
though ARIMA and the neural direct-output heads generate them differently.

Numerical convergence uses a uniform policy: L-BFGS with at most 200
iterations, followed only upon nonconvergence by Powell with at most 2,000
iterations for the same order and data. Optimizer choice does not use
validation or test forecast errors. Nonconverged candidates are logged and
excluded; no different forecasting model is silently substituted. This
policy was introduced after a pre-2024 refit failed in the initial
numerical run; that attempt is retained separately. ARIMA has no seed
ensemble or artificial across-seed standard deviation. Architecture sizes
and search procedures are not equalized, so this benchmark compares the
declared configurations rather than proving an intrinsic family ranking.
